\documentclass[11pt,reqno]{amsart}

\usepackage[margin=1in]{geometry}
\usepackage{amsmath,amssymb,amsthm}
\usepackage{booktabs}
\usepackage{microtype}
\usepackage[colorlinks=true,linkcolor=blue,citecolor=blue]{hyperref}

\newtheorem{theorem}{Theorem}[section]
\newtheorem{lemma}[theorem]{Lemma}
\newtheorem{proposition}[theorem]{Proposition}
\newtheorem{corollary}[theorem]{Corollary}
\theoremstyle{definition}
\newtheorem{definition}[theorem]{Definition}
\newtheorem{remark}[theorem]{Remark}

\newcommand{\R}{\mathbb{R}}
\newcommand{\T}{\mathbb{T}}
\newcommand{\Z}{\mathbb{Z}}
\newcommand{\uvec}{\mathbf{u}}
\newcommand{\vvec}{\mathbf{v}}
\newcommand{\wvec}{\boldsymbol{\omega}}
\newcommand{\divg}{\operatorname{div}}
\newcommand{\curl}{\operatorname{curl}}

\numberwithin{equation}{section}

\title[Metriplectic Phase-Space Completion for 3D Navier--Stokes]{Geometric Phase-Space Completion and Conditional Global Regularity for the 3D Incompressible Navier--Stokes Equations}

\author{Benjamin Stanley Frohman}
\date{\today}

\begin{document}

\begin{abstract}
Let \(\uvec_0\in C^\infty(\T^3;\R^3)\) be a smooth, divergence-free initial velocity field on the flat three-dimensional torus \(\T^3=\R^3/(2\pi\Z)^3\). We formulate a geometric completion of Arnold's Lie--Poisson framework on \(\mathrm{SDiff}(\T^3)\) via a three-legged metriplectic architecture: (1) the skew-symmetric Arnold Lie--Poisson bracket governing ideal transport, (2) a symmetric metriplectic dissipation sector incorporating viscous damping alongside a degenerate projected metric correction \(B(F,G)\) (the Frohmanian Symplectic Tether), and (3) Marsden--Weinstein--Ratiu coadjoint orbit reduction. Because \(B(F,H)\equiv 0\) for the kinetic-energy Hamiltonian \(H\), the generated dynamics recover the exact, unmodified 3D Navier--Stokes vorticity transport equation. Within this framework we record three rigorous regularity results: a sharp conditional global continuation theorem via critical \(L^1_t L^\infty_x\) strain integrability (Theorem~A), unconditional local-in-time smoothness on the deterministic Kato lifespan (Theorem~B), and unconditional global smoothness for small critical \(\dot{H}^{1/2}\) data (Theorem~C). Finally, we isolate the supercritical \(L^2_t L^2_x\to L^1_t L^\infty_x\) scaling gap as the precise open frontier for geometric vorticity-coherence analysis. No unconditional large-data claim is made; Clay Statement~(B) remains open.
\end{abstract}

\maketitle

\section{Statement of the Problem: The Clay Millennium Problem and the Supercriticality Barrier}

\subsection{Classical PDE Posing (Fefferman Formulation)}
Let \(\T^3=\R^3/(2\pi\Z)^3\) denote the flat three-dimensional torus. The 3D incompressible Navier--Stokes equations in velocity-pressure form are
\begin{equation}\label{eq:NSE_vel}
\begin{cases}
\partial_t\uvec+(\uvec\cdot\nabla)\uvec+\nabla p=\nu\Delta\uvec,\\
\divg\uvec=0,\\
\uvec(x,0)=\uvec_0(x),
\end{cases}
\end{equation}
where \(\uvec_0\in C^\infty(\T^3;\R^3)\) is smooth and divergence-free, \(p\) is the scalar pressure, and \(\nu>0\) is the kinematic viscosity. Taking the curl of~\eqref{eq:NSE_vel} eliminates the pressure gradient and yields the equivalent vorticity transport equation for \(\wvec=\curl\uvec\):
\begin{equation}\label{eq:NSE_vort}
\partial_t\wvec+(\uvec\cdot\nabla)\wvec=(\wvec\cdot\nabla)\uvec+\nu\Delta\wvec,\qquad\divg\wvec=0.
\end{equation}
The Clay Millennium Prize Problem (Statement~B) asks whether, for every smooth divergence-free \(\uvec_0\), there exists a unique global classical solution \(\uvec\in C^\infty(\T^3\times[0,\infty);\R^3)\) (\(T^*=\infty\)), or whether a finite-time singularity (\(T^*<\infty\)) can form~\cite{Fefferman2000}.

\subsection{The Analytical Supercriticality Obstruction}
Testing~\eqref{eq:NSE_vel} against \(\uvec\) yields the unconditional global energy inequality, providing an \(L^\infty_t L^2_x\) bound on velocity and an \(L^2_t L^2_x\) bound on vorticity:
\begin{equation}\label{eq:energy_ineq}
\|\uvec(t)\|_{L^2(\T^3)}^2+2\nu\int_0^t\|\wvec(\tau)\|_{L^2(\T^3)}^2\,d\tau\le\|\uvec_0\|_{L^2(\T^3)}^2.
\end{equation}
Under the natural parabolic rescaling \(\uvec_\lambda(x,t)=\lambda\uvec(\lambda x,\lambda^2 t)\) on \(\R^3\), the kinetic energy scales as
\[
\|\uvec_\lambda(t)\|_{L^2(\R^3)}^2=\lambda^{-1}\|\uvec(\lambda^2 t)\|_{L^2(\R^3)}^2\to 0\quad\text{as }\lambda\to\infty.
\]
Because the scaling exponent is negative (\(-1<0\)), the \(L^2\) energy bound is \emph{supercritical}. Preventing finite-time blow-up requires controlling a scale-invariant (\(\lambda^0\)) critical norm, such as the \(L^1_t L^\infty_x\) integrability of the symmetric strain tensor \(S=\frac12(\nabla\uvec+(\nabla\uvec)^T)\), leaving a fundamental analytical gap between available \(L^2_t L^2_x\) bounds and required \(L^1_t L^\infty_x\) control.

\section{Geometric Completion of Phase Space}

Rather than modifying the underlying PDE, we embed~\eqref{eq:NSE_vort} into a geometrically complete infinite-dimensional phase space over the group of volume-preserving diffeomorphisms \(\mathrm{SDiff}(\T^3)\) with Lie algebra \(\mathfrak{g}=\mathfrak{sdiff}(\T^3)\).

\subsection{Leg 1: Reversible Lie--Poisson Sector}
On the dual space \(\mathfrak{g}^*\), ideal inviscid fluid motion is governed by Arnold's skew-symmetric Lie--Poisson bracket~\cite{Arnold1966}, written here in the paper's vorticity identification
\begin{equation}\label{eq:LP_bracket}
\{F,G\}_{\mathrm{LP}}=\int_{\T^3}\wvec\cdot\Bigl(\frac{\delta F}{\delta\wvec}\times\frac{\delta G}{\delta\wvec}\Bigr)\,dV.
\end{equation}
\begin{remark}[Jacobi identity by coadjoint reduction]\label{rem:jacobi}
We do not derive the Jacobi identity for~\eqref{eq:LP_bracket} from scratch. The bracket is the canonical reduced Poisson structure on \(\mathfrak{sdiff}(\T^3)^*\), obtained by infinite-dimensional Marsden--Weinstein reduction of the geodesic (kinetic-energy) Hamiltonian flow on \(T^*\mathrm{SDiff}(\T^3)\). Consequently the Jacobi identity holds as a theorem of Lie--Poisson / coadjoint-orbit reduction~\cite{Arnold1966,MW1983,AK1998}. (An equivalent curl-form presentation \(\int\wvec\cdot\bigl((\nabla\times\delta F/\delta\wvec)\times(\nabla\times\delta G/\delta\wvec)\bigr)\,dV\) appears under the usual identification of variational derivatives with divergence-free vector fields; we retain~\eqref{eq:LP_bracket} as the paper's written convention.)
\end{remark}
Equipped with the kinetic-energy Hamiltonian
\[
H[\wvec]=\frac12\int_{\T^3}|\uvec|^2\,dV,\qquad\uvec=K_{\mathrm{BS}}*\wvec,
\]
where \(K_{\mathrm{BS}}\) denotes the Biot--Savart kernel on \(\T^3\), one has \(\nabla\times(\delta H/\delta\wvec)=\uvec\) (equivalently \(\delta H/\delta\wvec=\boldsymbol{\psi}\) with \(\curl\boldsymbol{\psi}=\uvec\)). The bracket \(\{F,H\}_{\mathrm{LP}}\) generates the exact inviscid Euler dynamics tangent to coadjoint orbits \(\mathcal{O}\subset\mathfrak{g}^*\).

\subsection{Leg 2: Symmetric Metriplectic Sector and the Tether Correction}
To account for viscous dissipation (\(\nu>0\)) while preserving exact algebraic symmetry, we equip \(\mathfrak{g}^*\) with a symmetric, negative-semidefinite metriplectic bracket
\[
[F,G]_{\mathrm{total}}=[F,G]_{\mathrm{diss}}+B(F,G),
\]
defined by~\cite{Morrison1986}
\begin{align}
[F,G]_{\mathrm{diss}}&=-\nu\int_{\T^3}\Bigl(\nabla\times\frac{\delta F}{\delta\wvec}\Bigr)\cdot\Bigl(\nabla\times\frac{\delta G}{\delta\wvec}\Bigr)\,dV,\label{eq:diss_bracket}\\
B(F,G)&=-\kappa\int_{\T^3}|\wvec|^2\Bigl(\Pi_{\uvec}\Bigl(\nabla\times\frac{\delta F}{\delta\wvec}\Bigr)\cdot\Pi_{\uvec}\Bigl(\nabla\times\frac{\delta G}{\delta\wvec}\Bigr)\Bigr)\,dV,\label{eq:tether_bracket}
\end{align}
where \(\kappa=C_{\mathrm{CZ}}(3)>0\) is the 3D Calder\'on--Zygmund constant and \(\Pi_{\uvec}\) is the \(L^2\)-orthogonal projection onto \(\mathrm{span}\{\uvec\}^\perp\):
\[
\Pi_{\uvec}v=v-\frac{\langle v,\uvec\rangle_{L^2}}{\|\uvec\|_{L^2}^2}\uvec\quad(\uvec\neq 0).
\]
Separating \(\{F,G\}_{\mathrm{LP}}\) (skew-symmetric) from \([F,G]_{\mathrm{total}}\) (symmetric) leaves the Jacobi identity of the Poisson sector intact (Remark~\ref{rem:jacobi}) while ensuring that \(B(F,G)\) satisfies strict degeneracy with respect to the kinetic-energy Hamiltonian \(H\):
\begin{equation}\label{eq:degeneracy}
\nabla\times\Bigl(\frac{\delta H}{\delta\wvec}\Bigr)=\uvec\implies\Pi_{\uvec}\Bigl(\nabla\times\frac{\delta H}{\delta\wvec}\Bigr)=\Pi_{\uvec}(\uvec)\equiv 0\implies B(F,H)\equiv 0\quad\forall F.
\end{equation}
Consequently, the complete metriplectic evolution
\begin{equation}\label{eq:metriplectic_flow}
\frac{dF}{dt}=\{F,H\}_{\mathrm{LP}}+[F,H]_{\mathrm{total}}=\{F,H\}_{\mathrm{LP}}+[F,H]_{\mathrm{diss}}
\end{equation}
reproduces the exact, unmodified 3D Navier--Stokes vorticity transport equation~\eqref{eq:NSE_vort} without introducing artificial forcing terms. The correction \(B\) is \emph{not} claimed to be Poisson.

\subsection{Leg 3: Marsden--Weinstein--Ratiu Reduction}
Because the volume form \(dV\) and the quadratic weight \(|\wvec|^2\,dV\) transform covariantly under the coadjoint action of \(\mathrm{SDiff}(\T^3)\), the reversible bracket~\eqref{eq:LP_bracket} descends to a canonical symplectic structure on each coadjoint orbit \(\mathcal{O}\) via Marsden--Weinstein--Ratiu reduction~\cite{MW1983,MWR1984,AK1998}, while \([F,H]_{\mathrm{diss}}\) precisely parameterizes the transverse dissipative drift across orbits driven by kinematic viscosity.

\section{Rigorous Regularity Theorems within the Framework}

Within the metriplectic formulation~\eqref{eq:metriplectic_flow}, the regularity theory of 3D Navier--Stokes resolves into three complete theorems.

\begin{theorem}[Conditional Global Regularity --- Form~B]\label{thm:A}
Let \(\uvec_0\in H^s(\T^3;\R^3)\) (\(s>5/2\)) satisfy \(\divg\uvec_0=0\), and let \(\uvec\in C([0,T^*);H^s(\T^3))\) be the maximal local Kato solution generated by~\eqref{eq:metriplectic_flow}. Assume that Form~B holds: for every finite \(T<T^*\), there exists a data-controlled constant \(C_*(\|\uvec_0\|_{H^s},\nu,T)<\infty\) such that
\begin{equation}\label{eq:FormB}
\int_0^T\|S(t)\|_{L^\infty(\T^3)}\,dt\le C_*(\|\uvec_0\|_{H^s},\nu,T).
\end{equation}
Then \(T^*=+\infty\) and \(\uvec\in C^\infty(\T^3\times(0,\infty);\R^3)\). Conversely, any globally smooth solution satisfies~\eqref{eq:FormB} on every finite interval \([0,T]\). In particular, the one-directional Beale--Kato--Majda continuation implication
\[
\int_0^T\|S(t)\|_{L^\infty}\,dt<\infty\implies T^*>T
\]
holds for every finite \(T\).
\end{theorem}

\begin{proof}
Decomposing \(\nabla\uvec=S+\Omega\) into symmetric and skew-symmetric parts, we have \((\wvec\cdot\nabla)\uvec=S\wvec\). Testing~\eqref{eq:NSE_vort} against \(|\wvec|^{p-2}\wvec\) and passing to the limit \(p\to\infty\) yields the differential inequality for \(M(t):=\|\wvec(t)\|_{L^\infty(\T^3)}\):
\[
\frac{d}{dt}\log M(t)\le\|S(t)\|_{L^\infty(\T^3)}.
\]
Integrating over \([0,T]\) and applying~\eqref{eq:FormB} bounds the vorticity supremum exponentially:
\[
M(T)\le M(0)\exp\Bigl(\int_0^T\|S(t)\|_{L^\infty(\T^3)}\,dt\Bigr)\le M(0)\,e^{C_*}.
\]
The sharp 3D Calder\'on--Zygmund logarithmic inequality~\cite{BKM1984,KozonoTaniuchi2000} states
\[
\|\nabla\uvec(t)\|_{L^\infty(\T^3)}\le C\Bigl(1+\|\wvec(t)\|_{L^\infty(\T^3)}\log\bigl(e+\|\uvec(t)\|_{H^s(\T^3)}\bigr)+\|\wvec(t)\|_{L^2(\T^3)}\Bigr).
\]
Substituting this bound into the standard Kato energy inequality
\[
\frac{d}{dt}\|\uvec(t)\|_{H^s}^2\le C_s\|\nabla\uvec(t)\|_{L^\infty}\|\uvec(t)\|_{H^s}^2
\]
and setting \(Y(t):=\log(e+\|\uvec(t)\|_{H^s}^2)\) produces the linear differential inequality
\[
\frac{d}{dt}Y(t)\le C'\bigl(1+\|\wvec(t)\|_{L^2(\T^3)}+M(t)\bigr)Y(t).
\]
By Gronwall's lemma, since \(\int_0^T\|\wvec(t)\|_{L^2}\,dt\) is bounded by~\eqref{eq:energy_ineq} and \(M(t)\) is uniformly bounded on \([0,T]\), the norm \(\|\uvec(t)\|_{H^s}\) cannot blow up as \(t\uparrow T^*\). By the Beale--Kato--Majda continuation principle, \(T^*=+\infty\). Parabolic regularity for the heat operator \((\partial_t-\nu\Delta)\) bootstraps \(\uvec\) to \(C^\infty(\T^3\times(0,\infty))\). The converse follows from continuous Sobolev embedding \(H^s(\T^3)\hookrightarrow W^{1,\infty}(\T^3)\) for \(s>5/2\).
\end{proof}

\begin{theorem}[Unconditional Local-in-Time Regularity --- Kato Lifespan]\label{thm:B}
For arbitrary large divergence-free initial data \(\uvec_0\in H^s(\T^3;\R^3)\) (\(s>5/2\)), there exists a deterministic lifespan
\[
T_{\mathrm{loc}}\ge\frac{1}{C_s\|\uvec_0\|_{H^s(\T^3)}}
\]
on which Form~B holds unconditionally and \(\uvec\in C^\infty(\T^3\times(0,T_{\mathrm{loc}}))\).
\end{theorem}

\begin{proof}
We record the three standard pillars of the Kato theory~\cite{Kato1972}.

\paragraph{(a) Full Kato--Ponce commutator bound.}
Applying \(\Lambda^s=(I-\Delta)^{s/2}\) to~\eqref{eq:NSE_vel} and taking the \(L^2\) pairing with \(\Lambda^s\uvec\) produces, after the usual cancellation of the pressure and the divergence-free constraint,
\[
\frac12\frac{d}{dt}\|\uvec\|_{H^s}^2+\nu\|\nabla\uvec\|_{H^s}^2
\le C_s\Bigl(\|\nabla\uvec\|_{L^\infty}\|\uvec\|_{H^s}+\|\nabla\uvec\|_{L^\infty}\|\nabla\uvec\|_{H^{s-1}}\Bigr)\|\uvec\|_{H^s}.
\]
For \(s>5/2\) the second summand is absorbed into the first by the Sobolev embedding \(H^{s-1}\hookrightarrow L^\infty\) on the gradient factor (or by interpolating \(\|\nabla\uvec\|_{H^{s-1}}\le C\|\uvec\|_{H^s}\)), yielding the closed differential inequality
\begin{equation}\label{eq:kato_riccati}
\frac{d}{dt}\|\uvec\|_{H^s}\le C_s\|\uvec\|_{H^s}^2.
\end{equation}

\paragraph{(b) Existence via mollified/Galerkin approximations.}
Let \(\uvec_0^\varepsilon=J_\varepsilon\uvec_0\) be a family of smooth mollifications (or finite-mode Galerkin projections) of the initial datum, with \(\|\uvec_0^\varepsilon\|_{H^s}\le\|\uvec_0\|_{H^s}\). The corresponding approximate solutions \(\uvec^\varepsilon\) satisfy~\eqref{eq:kato_riccati} with constants independent of \(\varepsilon\). Integrating the Riccati inequality on the common interval
\[
\bigl[0,T_{\mathrm{loc}}\bigr),\qquad T_{\mathrm{loc}}=\frac{1}{2C_s\|\uvec_0\|_{H^s}},
\]
gives a uniform bound \(\sup_{t<T_{\mathrm{loc}}}\|\uvec^\varepsilon(t)\|_{H^s}\le 2\|\uvec_0\|_{H^s}\). An Aubin--Lions compactness argument (boundedness in \(L^\infty_t H^s\) and of \(\partial_t\uvec^\varepsilon\) in a negative Sobolev space) produces a limit \(\uvec\in L^\infty([0,T_{\mathrm{loc}});H^s)\) solving~\eqref{eq:NSE_vel} distributionally; continuity in time with values in \(H^s\) follows by the standard energy-identity argument for strong solutions. Parabolic bootstrap then upgrades to \(\uvec\in C^\infty(\T^3\times(0,T_{\mathrm{loc}}))\).

\paragraph{(c) \(L^2\) uniqueness.}
If \(\uvec\) and \(\vvec\) are two strong solutions on \([0,T]\) with the same initial datum, set \(\mathbf{w}=\uvec-\vvec\). Testing the difference of the momentum equations against \(\mathbf{w}\) yields
\[
\frac12\frac{d}{dt}\|\mathbf{w}\|_{L^2}^2+\nu\|\nabla\mathbf{w}\|_{L^2}^2
=-\int_{\T^3}\bigl((\mathbf{w}\cdot\nabla)\vvec\bigr)\cdot\mathbf{w}\,dV
\le\|\nabla\vvec\|_{L^\infty}\|\mathbf{w}\|_{L^2}^2.
\]
Gronwall's lemma with \(\mathbf{w}(0)=0\) forces \(\mathbf{w}\equiv 0\). (The gradient factor is that of the second solution \(\vvec\), consistent with the subtraction \(\uvec-\vvec\).)

Via Sobolev embedding, Form~B holds on every compact subinterval of \([0,T_{\mathrm{loc}})\) with
\[
C_*(\|\uvec_0\|_{H^s},\nu,T)=-\frac{C_{\mathrm{emb}}}{C_s}\log\bigl(1-C_s T\|\uvec_0\|_{H^s}\bigr)<\infty.
\]
\end{proof}

\begin{theorem}[Unconditional Global Regularity for Small Critical Data --- Fujita--Kato]\label{thm:C}
There exists an absolute dimensionless constant \(\varepsilon>0\) such that whenever \(\uvec_0\in H^s(\T^3;\R^3)\) (\(s>5/2\)) satisfies the scale-invariant smallness condition
\[
\|\uvec_0\|_{\dot{H}^{1/2}(\T^3)}\le\varepsilon\nu,
\]
Form~B holds globally on \([0,\infty)\), yielding \(T^*=+\infty\) and \(\uvec\in C^\infty(\T^3\times(0,\infty))\).
\end{theorem}

\begin{proof}
Testing~\eqref{eq:NSE_vel} against \((-\Delta)^{1/2}\uvec\) and applying product estimates in homogeneous Sobolev spaces~\cite{FujitaKato1964} gives
\[
\frac12\frac{d}{dt}\|\uvec(t)\|_{\dot{H}^{1/2}}^2+\nu\|\uvec(t)\|_{\dot{H}^{3/2}}^2\le C_{\mathrm{FK}}\|\uvec(t)\|_{\dot{H}^{1/2}}\|\uvec(t)\|_{\dot{H}^{3/2}}^2.
\]
Choosing \(\varepsilon<1/(2C_{\mathrm{FK}})\) ensures that \((\nu-C_{\mathrm{FK}}\|\uvec_0\|_{\dot{H}^{1/2}})>\nu/2>0\) at \(t=0\), trapping \(\|\uvec(t)\|_{\dot{H}^{1/2}}\le\|\uvec_0\|_{\dot{H}^{1/2}}\) for all \(t\ge 0\) and forcing
\[
\int_0^\infty\|\uvec(t)\|_{\dot{H}^{3/2}}^2\,dt\le\frac1\nu\|\uvec_0\|_{\dot{H}^{1/2}}^2.
\]
Interpolating this global \(\dot{H}^{3/2}\) dissipation into higher-order \(H^s\) norms yields control of \(\|S(t)\|_{L^\infty}\) in \(L^1([0,\infty))\) by \(\|\uvec_0\|_{H^s}\) and \(\nu\). Theorem~\ref{thm:A} then yields global smoothness for \(t>0\).
\end{proof}

\section{Status Matrix, Lean Scope, and Open Frontier}

\begin{table}[htbp]
\centering
\caption{Status matrix for the results recorded in this manuscript.}
\label{tab:status}
\begin{tabular}{@{}llp{7.2cm}@{}}
\toprule
\textbf{Item} & \textbf{Status} & \textbf{Content} \\
\midrule
Theorem~A & Conditional & Form~B \(\Rightarrow\) global smoothness (\(t>0\)); converse on finite horizons \\
Theorem~B & Unconditional Local & Kato lifespan; smoothness for \(t>0\) \\
Theorem~C & Unconditional Small-Data Global & \(\|\uvec_0\|_{\dot{H}^{1/2}}\le\varepsilon\nu\) \\
Metriplectic geometry & Exact / Canonical & \(\{F,G\}_{\mathrm{LP}}\) skew; \([F,G]_{\mathrm{diss}}+B\) symmetric; \(B(F,H)\equiv 0\) \\
BKM continuation & One-directional & \(\int_0^T\|S\|_{L^\infty}\,dt<\infty\implies T^*>T\) \\
Lean formalization & Machine-checked (conditional) & Metriplectic algebra, degeneracy, alignment desingularization, Gronwall logic over abstract PDE input types; \(0\) \texttt{sorry}; axioms \texttt{[propext, Classical.choice, Quot.sound]}; Lean~/~Mathlib \texttt{v4.30.0-rc1} \\
Clay Statement~(B) & Open & The \(L^2_t L^2_x\to L^1_t L^\infty_x\) scaling gap \\
\bottomrule
\end{tabular}
\end{table}

The companion Lean module formalizes the metriplectic algebra, the degeneracy identity \(B(F,H)\equiv 0\), the alignment-kernel / Constantin--Fefferman desingularization bounds, and the integrating-factor Gronwall logic over abstract solution-data interfaces. Analytic Navier--Stokes facts (Clay Statement~(B), the BKM criterion, smooth-implies-bounded-strain, and the differential inequality) remain assumed structure fields; the Lean development does not claim an unconditional resolution of the Millennium problem.

\subsection{Isolation of the \(L^2_t L^2_x\to L^1_t L^\infty_x\) Scaling Gap}
As demonstrated by Theorems~\ref{thm:A}--\ref{thm:C}, the large-data global regularity problem reduces to establishing a data-only \emph{a priori} bound on \(\int_0^T\|S(t)\|_{L^\infty}\,dt\) (Form~B). The phase-space geometry makes transparent why standard energy methods cannot bridge this gap:
\begin{enumerate}
\item \textbf{Supercritical Dissipation vs.\ Critical Strain.} The only coercive global bound generated unconditionally by \([F,H]_{\mathrm{diss}}\) is the \(L^2_t L^2_x\) enstrophy dissipation~\eqref{eq:energy_ineq}, which scales as \(\lambda^{-1}\) under parabolic rescaling. Form~B requires \(L^1_t L^\infty_x\) control (\(\lambda^0\)).
\item \textbf{Necessity of Viscous Coupling.} Because \(B(F,H)\equiv 0\) leaves the reversible Euler sector untouched---and inviscid 3D Euler admits finite-time singularities in low-regularity classes---no purely symplectic orbit invariant can close Form~B without exploiting viscous dissipation.
\end{enumerate}

\subsection{Directions for Geometric-Coherence Analysis}
While the current degenerate bracket \(B(F,G)\) preserves the unmodified flow without injecting artificial damping, its projected structure \(\Pi_{\uvec}\) points toward geometric regularity criteria. Expressing the strain via the Constantin--Fefferman singular integral kernel~\cite{ConstantinFefferman1993} shows that vortex stretching cancels when the vorticity direction \(\boldsymbol{\xi}=\wvec/|\wvec|\) remains Lipschitz-coherent in intense regions. Whether localized coherence bounds on \(\boldsymbol{\xi}\) can be coupled to \([F,G]_{\mathrm{diss}}\) to bridge the scaling gap for arbitrary large data remains open.

\bibliographystyle{amsplain}
\bibliography{references}

\end{document}
